\documentclass[11pt,a4paper]{article}
\usepackage{amsmath,amssymb,amsthm}
\usepackage[margin=1.05in]{geometry}
\usepackage{longtable}
\usepackage{hyperref}

\theoremstyle{plain}
\newtheorem{theorem}{Theorem}
\newtheorem{conjecture}[theorem]{Conjecture}
\newtheorem{corollary}[theorem]{Corollary}
\newtheorem{proposition}[theorem]{Proposition}
\newtheorem{lemma}[theorem]{Lemma}
\theoremstyle{definition}
\newtheorem{remark}[theorem]{Remark}

\renewcommand{\SS}{\mathfrak{S}}
\newcommand{\AAA}{\mathfrak{A}}

\title{A negative map of the Elliott--Halberstam route to binary
Goldbach,\\ and the wall's own law}

\author{[author]}
\date{Version 1\\
\small All measurements reproducible from the accompanying repository}

\begin{document}
\maketitle

\begin{abstract}
Huang and Li proved that binary Goldbach for large even $N$ follows
from $EH_\mu(N^{\theta'})$ for a single $\theta' > 1/2$. This paper
maps that hypothesis from both of its sides and then studies the
single scalar it reduces to.

On the \emph{demand} side we prove one unconditional theorem
(Theorem~\ref{thm:A}): the M\"obius-weighted, fixed-class correlation
sum with weight $w_k = 1$ is $\ll_A N(\log N)^{-A}$, so Huang--Li's
$E_4$ consumption of $EH_\mu$ is unnecessary and the whole demand
collapses to the single scalar $E_3$. We then show
(Theorem~\ref{thm:C}) that this scalar is \emph{equivalent} to binary
Goldbach, by an unconditional identity which is Huang--Li's own
equation (22); the root cause is $\mu * \log = \Lambda$. The
demand-side search for a weaker sufficient condition is therefore
closed at the level of identities. Net progress toward Goldbach: zero.

On the \emph{supply} side we state one formal conjecture
(Conjecture~\ref{conj:L}): every M\"obius family probed factorizes as
a deterministic local mask times an exactly Gaussian fluctuation, so
that every ``sub-random'' reading in the literature of this problem
that we could reproduce is mask accounting. We report eighteen
pre-registered closures --- five route adjudications against the
source papers' explicit lemma hypotheses, ten kill-tested technique
designs, and three representation-class experiments --- together with
the five structural constraints they impose on any future technique.

The demand side reduces everything to
$C(N)=\sum_{n<N}\Lambda(n)\mu(N-n) = o(N)$, and
Sections~\ref{sec:wall}--\ref{sec:floor} give that quantity its own
law. We prove its exact second moment in closed form
(Proposition~\ref{prop:V}), identify the excess over random signs as a
prime-pair-weighted Chowla correlation (Proposition~\ref{prop:W}),
give an exact identity for the aggregate second moment
(Lemma~\ref{lem:MP}), and measure the deterministic part --- the
\emph{location mask} --- against an exactly computed floor
(Lemma~\ref{lem:cellmom}). The mask's shape is the singular series of
the shift; its size is accounted for to within six percent; its decay
is not explained by that mechanism, and we show why not
(Proposition~\ref{prop:scaleinv}). Along the way we record that the
uncertainty of a cell mean falls like $(\log N)^{-1/2}$ and not like
$n_c^{-1/2}$ (Proposition~\ref{prop:coh}), which changes the width of
every interval in this subject built from a count.

The deliverable is the map and the law, not a theorem toward Goldbach.
We state throughout what is and is not claimed.
\end{abstract}

\tableofcontents

\section{Introduction}

\subsection{The conditional frame}

Let $N$ be a large even integer,
$\tilde r(N) = \sum_{n<N}\Lambda(n)\Lambda(N-n)$, and $\SS(N)$ the
singular series. Write $EH_\mu(Q)$ for the assertion that for every
$A>0$
\[
  \sum_{q\le Q}\max_{y<N}\max_{(a,q)=1}
  \Bigl|\sum_{\substack{n\le y\\ n\equiv a\,(q)}}\Lambda(n)\mu(N-n)
   -\frac{1}{\varphi(q)}\sum_{n\le y}\Lambda(n)\mu(N-n)\Bigr|
  \ll_A \frac{N}{(\log N)^A}.
\]
Huang and Li~\cite{HL}, continuing Pan~\cite{Pan}, proved that
$EH(N^{\theta}(\log N)^{2A+8})$ together with $EH_\mu(N^{1-\theta})$
implies $\tilde r(N)\ge \SS(N)(1-A(N))N + O(N(\log N)^{-A})$, and
deduced (their Corollary~1) that in view of Bombieri--Vinogradov,
binary Goldbach for all sufficiently large even integers follows from
$EH_\mu(N^{\theta'})$ alone for any single $\theta' > 1/2$.

This is the sharpest conditional frame we are aware of: one sentence
stands between classical technology and Goldbach. This paper asks what
that sentence actually costs.

\subsection{The organizing principle: two couplings}

The hypothesis contains a product of two M\"obius factors, and there
are exactly two ways their arguments can be coupled. This turns out to
organize everything.

\begin{itemize}
\item \textbf{Divisibility-coupled} (the \emph{demand} side). Inside
  Huang--Li's proof, $EH_\mu$ is consumed only in the fixed residue
  class $n \equiv N \pmod k$, i.e.\ $k \mid N-n$: one M\"obius
  argument divides the other. Section~\ref{sec:demand}.
\item \textbf{Difference-coupled} (the \emph{supply} side). Any attempt
  to \emph{prove} $EH_\mu$ by dispersion arrives at the dilate field
  $D(k) = \sum_m \mu(m)\mu(N-mk)$, where the two arguments are related
  by a difference, not a divisibility. Section~\ref{sec:supply}.
\end{itemize}

\noindent
The two behave completely differently, and each is closed for its own
reason:

\begin{center}
\begin{tabular}{p{0.20\textwidth}p{0.34\textwidth}p{0.34\textwidth}}
\hline
 & \textbf{Divisibility-coupled} & \textbf{Difference-coupled}\\
\hline
Divisor switch & exact, and free & exact, but useless\\
M\"obius lands on & the \emph{short} variable, $m\le N^{1/2-\delta}$ &
 the \emph{long} variable, $m \asymp N/K \ge N^{2/3}$\\
Consequence & Bombieri--Vinogradov closes it: Theorem~\ref{thm:A} &
 no known machine applies\\
But & the weighted version is \emph{equivalent} to Goldbach
 (Theorem~\ref{thm:C}) & no named coupling surface found, internal or
 external (\S\ref{sec:closures})\\
\hline
\end{tabular}
\end{center}

\subsection{What this paper claims}

\begin{itemize}
\item \textbf{Claimed}: Theorem~\ref{thm:A}, Corollary~\ref{cor:B},
  Theorem~\ref{thm:C}, Theorems~\ref{thm:D} and~\ref{thm:Dprime},
  Propositions~\ref{prop:E}, \ref{prop:Dpp}, \ref{prop:V},
  \ref{prop:W}, \ref{prop:scaleinv} and~\ref{prop:coh},
  Lemmas~\ref{lem:MP}, \ref{lem:cellmom}, \ref{lem:coin}
  and~\ref{lem:placebo}; the defect in \cite{HL}'s equation (18) and
  its repair; all measurements, which are reproducible; the closure
  table of \S\ref{sec:closures}.
\item \textbf{Conjectured}: Conjectures~\ref{conj:L}
  and~\ref{conj:wall}.
\item \textbf{Not claimed}: any theorem toward Goldbach.
  Theorem~\ref{thm:A} is unconditional but Goldbach-neutral: it
  removes exactly the half of the demand that carries no Goldbach
  content.
\end{itemize}

\section{The demand side}\label{sec:demand}

\subsection{Where $EH_\mu$ is actually consumed}

For $(k,N)=1$ and $1\le t<N$ put
\[
  E_\mu(t;k) = \sum_{\substack{n\le t\\ n\equiv N\,(k)}}
    \Lambda(n)\mu(N-n) - \frac{1}{\varphi(k)}\sum_{n\le t}
    \Lambda(n)\mu(N-n).
\]
With $K = (N-1)/\alpha$, Huang--Li's two consumption sites are
\[
  E_3(\alpha)=\!\!\sum_{\substack{k<K\\(k,N)=1}}\!\!\mu(k)\log k\,
  E_\mu(N;k),
  \qquad
  E_4(\alpha)=\!\!\sum_{\substack{k<K\\(k,N)=1}}\!\!\mu(k)\,
  \bigl[\text{$E_\mu$ with the extra weight }\log(N-n)\bigr].
\]
They differ in exactly one respect: the weight $w_k$ is $\log k$ in
$E_3$ and $1$ in $E_4$ (the $\log(N-n)$ inside $E_4$ is
$k$-independent and comes off by partial summation). In both cases
Huang--Li discard the signs $\mu(k)$ on the first line by the triangle
inequality, and only then appeal to $EH_\mu$. Keeping the signs is
what the following results exploit.

\subsection{The theorem}

\begin{theorem}[$w_k = 1$]\label{thm:A}
Fix $\theta' \in (1/2,1)$ and put $K = N^{\theta'}$. Then for every
$A>0$,
\[
  \sup_{1\le t<N}\Bigl|
   \sum_{\substack{k<K\\(k,N)=1}}\mu(k)\,E_\mu(t;k)\Bigr|
  \;\ll_{A,\theta'}\; \frac{N}{(\log N)^A}.
\]
Every ingredient except Bombieri--Vinogradov contributes only
$O(Ne^{-c\sqrt{\log N}})$.
\end{theorem}

The mechanism is one line. The class $n \equiv N \pmod k$ is
$k \mid N-n$, so the $k$-sum is a divisor sum over $u = N-n$;
switching and writing $u = mk$ gives, for squarefree $u$,
$\mu(u)\mu(k) = \mu(m)\mu^2(k)$ --- the M\"obius on the long variable
$k$ squares itself away, and the surviving M\"obius sits on the short
variable $m < N^{1-\theta'} \le N^{1/2-\delta}$. That is the classical
``M\"obius on the short variable plus Bombieri--Vinogradov''
configuration. Full proofs are in the companion note~\cite{ThmA}.

The bound is $N(\log N)^{-A}$ for each fixed $A$ and no better. It is
not $Ne^{-c\sqrt{\log N}}$: Bombieri--Vinogradov's internal
Siegel--Walfisz range is $q\le(\log N)^B$ with $B$ fixed, and that is
what caps the saving.

\begin{corollary}\label{cor:B}
$E_4(\alpha) \ll_A N(\log N)^{-A}$ unconditionally, via the exact
identity $E_4(\alpha) = \int_1^{N-1} T_1(t)\,dt/(N-t)$. Hence
Huang--Li's Lemma 4 is not needed and the entire $EH_\mu$ demand
collapses to the single scalar $E_3(\alpha)$.
\end{corollary}

\begin{theorem}[$w_k = \log k$: the permanent closure]\label{thm:C}
Unconditionally,
\[
  E_3(\alpha) = \sum_{n<N}\Lambda(n)\Lambda(N-n)
   - \SS(N)\Bigl(N - \sum_{n<N}\Lambda(n)\mu(N-n)\Bigr)
   + O_A\!\left(N(\log N)^{-A}\right),
\]
which is Huang--Li's own equation (22). Hence
$E_3(\alpha)\ll_A N(\log N)^{-A}$ --- the weakest form of the
hypothesis their argument needs --- is \emph{equivalent} to
$\tilde r(N)\sim\SS(N)N$.
\end{theorem}

\begin{remark}
The two theorems are the same computation with two weights. For
$w_k = \log k$ the complete divisor sum is $-\Lambda$ rather than an
indicator, because $\mu * \log = \Lambda$ --- and that identity is
where Huang--Li \emph{start} (their (10)). The $\log k$ weight is the
carrier of the Goldbach content, so every divisor switch must hand it
back. This is closure at the level of identities, not of estimates: no
choice of $\theta'$, truncation, or smoothing can evade it.
\end{remark}

\subsection{The interior of the weight space is empty}

Theorems~\ref{thm:A} and~\ref{thm:C} are the two endpoints of a design
space: the weight $w_k$ attached to $\mu(k)$ in the demand functional.
It is natural to look for a weight between them --- complete divisor
sum still cheap, main-term coefficient still of size $1$ --- that would
extract $C(N)=\sum_{n<N}\Lambda(n)\mu(N-n)$ from
Bombieri--Vinogradov alone. There is none.

Write $b := \mu*w$, so $w_k=\sum_{d\mid k}b_d$; every weight has this
form. For squarefree $u$ one has $\sum_{k\mid u}\mu(k)w_k=\mu(u)b_u$,
so the switch identity reads
$B_w\,C(N)=\sum_{u<N}\Lambda(N-u)\mu^2(u)b_u-\mathcal R_w$ with
$B_w=\sum_{k<K,(k,N)=1}\mu(k)w_k/\varphi(k)$. Two exact facts then
oppose each other:

\begin{itemize}
\item \emph{Extraction}:
  $B_w=\sum_{d<K}b_d\,\frac{\mu(d)}{\varphi(d)}\rho_{dN}(K/d)$ with
  $\rho \ll e^{-c\sqrt{\log x}}$, so $B_w$ is controlled by the part of
  $b$ sitting \emph{at} the truncation point $K=N^{\theta'}$.
\item \emph{BV-accessibility}: expanding $w$ inside the residual gives
  moduli $md$ with $m<N^{1-\theta'}$, so $b$ may only be supported on
  $d\le N^{\theta'-1/2-\delta}$.
\end{itemize}

\begin{theorem}[no weight extracts $C(N)$]\label{thm:D}
If $b=\mu*w$ is supported in $[1,N^{\theta'-1/2-\delta}]$ then
$\|b\|_1/|B_w| \gg \exp(c\sqrt{(1/2+\delta)\log N})$, and the switch
identity yields at best
$|C(N)|\ll \exp(c\sqrt{\tfrac12\log N})\,N(\log N)^{-A}$ --- no saving
of any power of $\log N$.
\end{theorem}

\begin{theorem}[the no-go survives $EH$]\label{thm:Dprime}
If $\Lambda$ has level of distribution $\theta_E<1$, the same argument
gives $\|b\|_1/|B_w| \gg \exp(c\sqrt{(1-\theta_E)\log N})$. Closing the
gap requires $\theta_E=1$: equidistribution of $\Lambda$ to moduli of
size $N$ itself, where each progression holds $O(1)$ terms.
\end{theorem}

\begin{proposition}[the circle method has zero margin]\label{prop:E}
Writing $C(N)=\int_0^1 S_\Lambda S_\mu e(-N\alpha)d\alpha$, both
standard estimates lie at or above the trivial bound $\psi(N)\sim N$:
Cauchy--Schwarz gives $\sim(6/\pi^2)^{1/2}N(\log N)^{1/2}$, exceeding
it by a growing factor; and any bound
$\sup_\alpha|S_\mu|\cdot\|S_\Lambda\|_1$ is at least
$\|S_\mu\|_2\|S_\Lambda\|_1\gg N$ by Parseval. The measured margin
$N/(\sup|S_\mu|\,\|S_\Lambda\|_1)$ is $0.168,0.175,0.158,0.152$ at
$N=2^{14},\dots,2^{20}$ --- below $1$ and decaying.
\end{proposition}

Davenport's uniform bound $S_\mu\ll_A N(\log N)^{-A}$ is useless here:
it saves against $N$, while the pairing needs the scale $N^{1/2}$,
which Parseval forbids. This is the parity obstruction in
circle-method language.

Theorem~\ref{thm:D} assumes $b$ is supported low enough for BV, which
excludes the most natural family: $w_k=f(\log k)$ with $f$ polynomial.
There $b=\mu*\log^D=\Lambda_D$, the generalised von Mangoldt function,
so the complete part splits by the number of prime factors of $u$ into
a Goldbach-type piece ($r=1$) and Chen-type pieces ($r\ge2$).

\begin{proposition}[polynomial weights]\label{prop:Dpp}
$D=0$ gives $B_w\ll e^{-c\sqrt{\log K}}$. For a monomial $x^D$,
$D\ge1$, every term of the complete part is nonnegative
($\Lambda\ge0$, $\Lambda_D\ge0$), so it is $\asymp N(\log N)^{D-1}$
with a fixed sign. Cancelling across monomials requires tuning against
the asymptotics of the $r=1$ piece, which at $D=1$ is the binary
Goldbach sum --- circular. No polynomial weight makes the complete
part a classically known object.
\end{proposition}

Measured at $N=10^6,4\cdot10^6,1.6\cdot10^7$: the two pieces of
$\mathrm{CP}_2$ are the same order, with ratio $0.771,0.790,0.810$
drifting towards $1$, and $\mathrm{CP}_2/(N\log N) = 2.886, 2.949,
2.997$. The closure rests on nonnegativity, not on one piece
dominating. The one analytically canonical tuning,
$f(x)=x^2-2\gamma x$, which makes $b$ mean-zero by killing the pole of
$\zeta W$ at $s=1$, moves the complete part by about five percent:
mean-zero is not smallness, because it removes the average of $b$ and
not its correlation with $\Lambda(N-\cdot)$.

The loss exponent $\sqrt{(1/2)\log N}$ is the $\sqrt N$ barrier itself:
it is the gap between where a weight must live to see $C(N)$ and where
Bombieri--Vinogradov permits it to live. This is a no-go for one
precisely specified method --- divisor switching with BV as the only
input --- over that method's entire weight space; it is not an
obstruction to other methods. Its purpose is to close the design space
that Theorems~\ref{thm:A} and~\ref{thm:C} opened.

\subsection{A defect in the published paper}

Equation (18) of \cite{HL} replaces the $n$-dependent constraint
$k < (N-n)/\alpha$, present in the definition of $S_2(\alpha)$, by the
$n$-free $k < (N-1)/\alpha$. The omitted term is
\[
  \Delta = \sum_{2\le m\le\alpha}\mu(m)\log m
    \sum_{\substack{k<(N-1)/\alpha\\ (k,m)=1}}\mu^2(k)\Lambda(N-mk),
\]
whose trivial bound $\ll N(\log N)^2$ exceeds the target. It closes by
the same machinery --- M\"obius again on the short variable --- under
hypotheses already assumed: unconditionally by Bombieri--Vinogradov in
the Corollary-1 regime, and under the assumed $EH$ in the Theorem-1
regime. Theorem 1 and Corollary 1 of \cite{HL} stand as stated.

\section{The supply side}\label{sec:supply}

\subsection{The consumable}

The chain's supply-side consumable, distilled, is
\[
  \textbf{E1 (weak form).}\quad
  D(k) = \!\!\sum_{\sqrt N < m \le N/k}\!\! \mu(m)\mu(N-mk),
  \qquad
  \sum_{k\sim K}|D(k)|^2 \ll (\log N)^{-2A-2}\sum_{k\sim K} M_k^2,
\]
for $K \le N^{1/3}$ dyadic, where $M_k$ is the length of the $m$-range.
A fixed log-power saving is the currency; $o(1)$ and $\log\log$ savings
are not consumable.

The normalisation is $M_k^2$, i.e.\ the \emph{trivial} bound, not
$M_k$, the square-root scale: the wall is
$T_{II} = \sum_{k\sim K}b_kD(k)$ with $|b_k|\ll\log N$, needed at
$N(\log N)^{-A}$, and Cauchy--Schwarz in $k$ gives
$\sum_k|D(k)|^2 \ll N^2/(K(\log N)^{2A+2})$ with
$\sum_{k\sim K}M_k^2 \asymp N^2/K$. The distinction is load-bearing:
at the square-root normalisation the demand would be
$\sum|D|^2/\sum M_k \ll 8.8\cdot10^{-6}$, which our own measurements
refute directly ($0.305, 0.310, 0.319$ at $N=10^8$). At the correct
normalisation the target is cleared by square-root cancellation with
margin $(N/K)(\log N)^{-2A-2}\to\infty$, though that margin is
invisible at accessible $N$.

\begin{remark}[the mandatory comparison]\label{rem:scale}
A target must be checked against one's own measurement of the same
quantity before it is used to adjudicate anything. Writing a target at
the square-root scale when the chain consumes it at the trivial scale
inflates every demand by a power of $N$, and no verdict passed on a
budget basis survives it. We record this as a standing rule because
the same slip occurred twice in the course of this work, and both times
it was caught only by that comparison.
\end{remark}

\subsection{Conjecture L}

\begin{conjecture}[the factorization law]\label{conj:L}
Each M\"obius family probed in this program --- prime-indexed dilate
pairs $C_{k,k'} = \sum_{p\sim P}\mu(N-pk)\mu(N-pk')$, integer-indexed
dilates and their pairs, and M\"obius sums over thin progressions ---
factorizes as
\[
  \text{field} \;=\; \mathbf M \times \mathbf G,
\]
where $\mathbf M$ is the deterministic local mask, computed exactly by
finite modular enumeration from the $v_q$-data of $(N,k,k')$, and
$\mathbf G$ is, on the surviving support, a fluctuation that is
exactly Gaussian at half-normal scale: no mean field, no class
structure, kurtosis $3$, Wishart-consistent pair spectrum.
\end{conjecture}

The evidence is summarized in Table~\ref{tab:L}; full stamps and
one-shot reproduction are in the repository.

\begin{table}[t]
\centering\small
\begin{tabular}{p{0.30\textwidth}p{0.62\textwidth}}
\hline
Level & Stamp\\
\hline
Pair statistics & free class $0.97$--$1.02$, kurtosis $2.99$--$3.03$
 (two $N$)\\
Exact cells & every viable $(v_2,v_3)$ cell $0.99$--$1.06$\\
Mask, blind & $\mathrm{corr}(\text{pred},\text{obs}) = 1.0000$,
 amplitude error $1.5\%$; annihilation fractions predicted exactly at
 a fresh $N$\\
Pair matrix & $\lambda_{\max}$ at the Wishart null dead center
 ($z = -0.19$)\\
E1 itself at $10^9$ & band ratios $0.966 / 0.950 / 0.922$ over
 $K \sim 10^3/3\cdot10^3/10^4$, no growth\\
E1 definitive grid & 4 $N$ $\times$ 2 bands $\times$ 300 $k$ with
 bootstrap SE: mean $z = +0.02$, 7/8 cells within $1.25\sigma$, the
 suspect cell settled by near-census\\
\hline
\end{tabular}
\caption{Measured support for Conjecture~\ref{conj:L}.}
\label{tab:L}
\end{table}

Conjecture~\ref{conj:L} is \emph{stronger} than the chain needs: E1
requires only square-root cancellation on average. What is missing is
therefore not knowledge of the structure --- the mask is an algorithm
and $\mathbf G$ is featureless in every measurement --- but a proof
technique for the amplitude of a featureless object.

\section{The wall's own law}\label{sec:wall}

The chain of Section~\ref{sec:demand} reduces the whole problem to one
scalar,
\[
  C(N) \;=\; \sum_{n<N}\Lambda(n)\mu(N-n) \;=\; o(N).
\]
This section gives that scalar its exact second moment, an exact
identity for its aggregate second moment, and the conjecture that
organizes what is measured about it.

\subsection{The exact scale}

\begin{proposition}[The exact scale]\label{prop:V}
Let
\[
  V(N) \;=\; \sum_{v<N}\mu^2(v)\,\Lambda(N-v)^2,
  \qquad
  W(N) \;=\; \sum_{w<N}\Lambda(w)^2,
\]
and let
\[
  \AAA(N) \;=\; \prod_{q\,\nmid\,N}\Bigl(1-\frac{1}{q(q-1)}\Bigr).
\]
Then $V(N) = W(N)\,\AAA(N)\,(1+o(1))$, and consequently
$V(N) \sim \AAA(N)\,N\log N$.
\end{proposition}

\begin{proof}[Derivation]
$\Lambda(w)^2$ is supported on prime powers with weight $(\log p)^2$,
so with $w=N-v$,
\[
  V(N) \;=\; \sum_{p^k<N}(\log p)^2\mu^2(N-p^k)
        \;=\; \sum_{p<N}(\log p)^2\mu^2(N-p) + O(\sqrt N\log^2 N),
\]
a $(\log p)^2$-weighted count of \emph{squarefree shifted primes}. Its
local density at a prime $q$ is $1$ when $q\mid N$ --- there
$q^2\mid N-p$ forces $q\mid p$, hence $p=q$, a single term --- and
$1-1/\varphi(q^2) = 1-1/(q(q-1))$ when $q\nmid N$, since the bad class
$p\equiv N \pmod{q^2}$ is then a unit class. The statement is Mirsky's
theorem on squarefree values of shifted primes \cite{Mir49} together
with partial summation; it is recalled here, not claimed.
\end{proof}

\paragraph{The local factor is $\AAA$, not $\SS$.}
The singular series $\SS(N)$ is Hardy--Littlewood's and is correct for
the Goldbach \emph{count}; the right object for the \emph{variance} of
$C(N)$ is $\AAA(N)$. Both are products over $q\mid N$, which is why
the substitution is easy to make and hard to see. Rescaling each
candidate to the measured mean, so that only its shape in $N$ is
judged, the residual standard deviation over every even
$N\le 1.6\cdot10^7$ is $0.000323$ for $\AAA$ against $0.245235$ for
$\SS$ --- a factor of $760$. Dividing through by $W(N)$ removes the
analytic factor exactly, so no asymptotic for $\sum(\log p)^2$ is
assumed: the ratio $V(N)/W(N)$ against $\AAA(N)$ has mean $1.000000$
with standard deviation $0.000145$ in the top octave, and agrees to
five decimals in each of the six radical cells from $2\mid N$ through
$2\cdot3\cdot5\cdot7\cdot11\cdot13\mid N$.

\begin{remark}[$\AAA$ is unambiguous as a band summary]
The ratio $V/W$ gives the same value to six decimals whether summarised
as a ratio of means or as a mean of ratios, in every octave; so does
the squarefree density $Q(N)/N$. Not every ratio in this subject is
that well behaved --- see Remark~\ref{rem:rho} --- so we state it.
\end{remark}

\subsection{The aggregate second moment, exactly}

\begin{lemma}[the second-moment identity]\label{lem:MP}
For any $X$,
\[
  \sum_{N\le X} C(N)^2 \;=\; \sum_{|h|<X} M(h)\,P(h),
  \qquad
  M(h) = \sum_{v} \mu(v)\mu(v+h),
  \quad
  P(h) = \sum_{w} \Lambda(w)\Lambda(w+h),
\]
the inner sums being over the ranges that keep both arguments below
$X$.
\end{lemma}

\begin{proof}
$C(N)$ is the convolution $(\Lambda * \mu)(N)$ evaluated with the
second argument reflected; expanding $C(N)^2$ as a double sum over
$(n,v)$ and $(n',v')$ with $n+v = n'+v' = N$ and summing over $N$
leaves the single constraint $n-n' = v'-v =: h$. Collecting the
$\Lambda$-pairs and the $\mu$-pairs separately gives the two factors.
\end{proof}

The identity is exact and unconditional, and it makes the aggregate
size of the wall a statement about two shifted correlations, one of
which ($P$) is the Hardy--Littlewood prime-pair sum and the other of
which ($M$) is the binary Chowla sum. It is the aggregate counterpart
of Proposition~\ref{prop:W} below, which is pointwise.

\subsection{The law}

\begin{conjecture}[The wall's law]\label{conj:wall}
\[
  C(N) \;=\; m(N) \;+\; \sqrt{V(N)}\;G(N),
\]
with $m(N)$ a deterministic \emph{location mask} indexed by which small
primes divide $N$, $V(N)$ the exact second moment of
Proposition~\ref{prop:V}, and $G(N)$ Gaussian.
\end{conjecture}

Four measurements, all on every even $N\le1.6\cdot10^7$ with the mask
removed by finite modular enumeration.

\begin{enumerate}
\item \textbf{The bulk is Gaussian, under this scale and no other.}
  Excess kurtosis $-0.0005$ ($z=-0.3$) and
  $E|G|/\mathrm{sd}(G)$ short of $\sqrt{2/\pi}$ by $0.00018$
  ($z=-0.8$), on $6.3\cdot10^6$ values, removing cell \emph{means}
  alone. Under an $\SS N$-based scale the same data give excess
  kurtosis $+0.1704$ at $z=98$: a normaliser that manufactures a heavy
  tail.

\item \textbf{The mask is real at every scale measured.} Against the
  exactly computed floor of Lemma~\ref{lem:cellmom} it clears
  Bonferroni in every octave to $1.6\cdot10^7$, at $\max_c|z_c| = 8.4$
  in the top band. At $N\approx10^8$ a $300$-versus-$300$ comparison
  gives the deep arm $0.9476\pm0.0293$ against the shallow
  $0.7238\pm0.0245$, a gap of $+0.2238\pm0.0056$ --- forty standard
  errors, and the one figure here that stakes a sign as well as a size.
  Section~\ref{sec:floor} is devoted to it.

\item \textbf{The tail is Gaussian too}, which is the half that
  matters, since $C(N)=o(N)$ constrains every $N$ and a bulk can match
  to five decimals while the tail is heavy. $\max|G|$ tracks the Gumbel
  law with mean deviation $+0.54\pm0.45$ from $E[\max]$ over eight
  octaves; aggregate tail counts against the Gaussian expectation give
  ratios $0.999$ at $t=3$, $0.997$ at $t=4$, $0.878$ at $t=5$; and the
  extremes are attained at generic $N$, not at deep radicals, so the
  mask removal is not leaking into the tail.

\item \textbf{$G$ is Gaussian in distribution but not phase-random in
  $\log N$.} Regressing $G$ on
  $\cos(\gamma\log N),\ \sin(\gamma\log N)$ for the first ten ordinates
  of the zeta zeros gives $R^2 = 3.90\cdot10^{-3}$ against a
  $200$-surrogate maximum of $5.09\cdot10^{-6}$, and every ordinate
  individually at $z\ge23$. This does not contradict item~1: a small
  oscillatory component leaves kurtosis untouched, and \emph{Gaussian
  in distribution} and \emph{phase-random in $\log N$} are different
  statements. The $0.39\%$ is a floor, not the share --- the zeros do
  not stop at the tenth.
\end{enumerate}

\section{The wall against a coin}\label{sec:coin}

$V(N)$ is exactly the random-sign second moment: if the $\mu(v)$ on the
surviving support were independent signs, $\mathrm{Var}\,C$ would equal
$V$. So the single number measuring the excess is
$\rho(N) = \mathrm{Var}\,C(N)/V(N)$, which is $1$ under random signs.

\begin{proposition}[The excess is a Chowla correlation]\label{prop:W}
With $u=N-p$ and $h=p'-p$,
\[
  \rho-1 \;=\; \frac1V\sum_{h\neq0} c(h)\,S(h),
  \qquad
  c(h)=\sum_{p'-p=h}(\log p)(\log p'),
  \quad
  S(h)=\bigl\langle \mu(u)\mu(u-h)\bigr\rangle .
\]
\end{proposition}

$S(h)$ is the binary Chowla correlation. Chowla's conjecture gives
$S(h)=o(1)$ for each fixed $h$, and its averaged form over $h$ is a
theorem \cite{MRT15}. Under that input $\rho\to1$, and then the wall is
\emph{exactly} square-root: nature over-delivers by a power of $\log$
and no more.

\begin{remark}[$\rho$ names three quantities]\label{rem:rho}
$\mathrm{mean}(C^2)/\mathrm{mean}(V)$, $\mathrm{mean}(C^2/V)$ and
$(\pi/2)\,\mathrm{mean}(|C|/\sqrt V)^2$ coincide only if $V$ is
constant across the band and $C/\sqrt V$ is exactly half-normal, and
neither holds. They agree to $0.75\%$ at $N\approx1.4\cdot10^7$ and
differ by $10.3\%$ at $10^5$, ordered as written. Proposition~\ref{prop:W}
is a statement about a ratio of totals and so wants the first; the
half-normal arm of the reproduction stamp reports $0.810\pm0.018$ at
$N\approx10^8$, which is $0.841$ in the units of
Proposition~\ref{prop:W}. Every quotation of $\rho$ must say which.
\end{remark}

\paragraph{What is measured.} The M\"obius autocorrelation sits at
$1.051$--$1.068$ times the random-sign floor $\sqrt{0.32264\,(X-h)}$
--- not $\sqrt X$, since the sum sees only $n$ with both $n$ and $n+h$
squarefree --- stably across five decades of shift. Reconstructing
$\rho-1$ from Proposition~\ref{prop:W} gives $-0.0976$ against a
measured $-0.18$, a factor $0.54$; the sign is negative, as $\rho<1$
requires.

\paragraph{Where the wall leans.} Gross mass by shift:
$h<10^3$ carries $1.1\%$, $10^3$--$10^4$ carries $3.0\%$,
$10^4$--$10^5$ carries $23.1\%$, $10^5$--$10^6$ carries $48.9\%$, and
above $10^6$ carries $23.8\%$ (the ranges cancel, net $-2.2\cdot10^{13}$
against gross $4.4\cdot10^{13}$, so shares of the net would mislead).
Small shifts, where Chowla is hardest and the averaged theorem weakest,
carry almost nothing: \textbf{the wall leans on the range where that
theorem is strongest.}

\paragraph{Major arcs.} On the rationals $j/q$ with small $q$, the
M\"obius exponential sums are markedly smaller than a coin's ---
by a factor $8.40$ at $q=3$ and $15.16$ at $q=5$. This is classical
major-arc cancellation and it is the mechanism behind $\rho<1$. It is
invisible to an unweighted average over frequencies, because the
periodogram weights each rational by $|\hat\Lambda(j/q)|^2$ and $q=3$
outweighs $q=143$ by more than three orders of magnitude; a comparison
of two frequency summaries must state its weight before the numbers can
be set against each other.

\subsection{Two controls, and what they invalidate}

\begin{lemma}[the coin control]\label{lem:coin}
Let $\varepsilon(v)=\pm1$ be arbitrary signs on $\{v : \mu(v)\neq0\}$
and zero elsewhere. Then $\varepsilon^2 = \mu^2$ pointwise, so
$V(N)$ is unchanged, and the field
$C_\varepsilon(N)=\sum_v\varepsilon(v)\Lambda(N-v)$ has the same exact
second moment as $C(N)$ for every $N$. Consequently any estimator whose
output is reproduced when $\mu$ is replaced by $\varepsilon$ is not
measuring $\mu$.
\end{lemma}

\begin{lemma}[the placebo key]\label{lem:placebo}
Permuting the cell labels across $N$ --- assigning to each $N$ the
label of some other $N$, by a fixed permutation of the label set ---
preserves every cell size and leaves the field $Z(N)=C(N)/\sqrt{V(N)}$
byte-identical. Any cell-indexed statistic that survives the
permutation is a property of the cell sizes and not of the
correspondence between cells and arithmetic.
\end{lemma}

These are trivial and they are load-bearing. Lemma~\ref{lem:coin}
invalidates a natural and tempting measurement: fitting $\rho(N)$
against a de-trended trend across a band of $N$ to extract a rate.
De-trending across a band removes real variance, because nearby $N$
share the same $\mu(v)\Lambda(N-v)$ terms and their $C(N)$ are
positively correlated; and the resulting estimator gives the same
output on $\varepsilon$ as on $\mu$. Proposition~\ref{prop:W} is an
identity and is untouched, but no fitted rate for $\rho-1$ is claimed
here. \textbf{Whether $\rho\to1$ --- whether the wall is exactly
square-root --- is open, and with one realisation of $\mu$ it may not
be answerable at all.} Lemma~\ref{lem:placebo} is what confirms, in
Section~\ref{sec:floor}, that the cell floor is a property of the
arithmetic and not of the partition.

\section{The location mask}\label{sec:floor}

The mask $m(N)$ of Conjecture~\ref{conj:wall} is the one deterministic
object in the law. This section states what is known about it: its
floor is exactly computable, its shape is the singular series of the
shift, its cause is the cell-to-divisibility correspondence, its size
is accounted for to within six percent, and its decay is not accounted
for by that mechanism at all.

Throughout, cells are indexed by \emph{depth} $d$, the number of
$3,5,7,11,13$ dividing $N$, and $Z(N)=C(N)/\sqrt{V(N)}$.

\subsection{The floor, in closed form}

\begin{lemma}[exact cell moments]\label{lem:cellmom}
Let a band of even $N$ be partitioned into cells, let $c$ be a cell of
size $n_c$ inside a band of size $n$, let $m_c$ and $\overline m$ be
the means of $Z$ over $c$ and over the band, and set
\[
  u_c(v) \;=\; \sum_{N\in c}\frac{\Lambda(N-v)}{\sqrt{V(N)}},
  \qquad
  Q_{cd} \;=\; \sum_v \mu^2(v)\,u_c(v)\,u_d(v).
\]
Then, when the $\mu(v)$ on the surviving support are replaced by
independent signs,
\[
  \mathrm{Var}\bigl(m_c-\overline m\bigr)
  \;=\; \frac{Q_{cc}}{n_c^{2}}
  \;-\; \frac{2\,Q_{ca}}{n_c\,n}
  \;+\; \frac{Q_{aa}}{n^{2}},
\]
where $a$ denotes the whole band. Every term is computable exactly by
one convolution; no simulation is involved.
\end{lemma}

\begin{proof}
Under independent signs $E[\varepsilon(v)\varepsilon(v')]
=\mu^2(v)\delta_{vv'}$, so
$\mathrm{Cov}(n_c m_c, n_d m_d) = Q_{cd}$ exactly. Expanding
$\mathrm{Var}(m_c-\overline m)$ bilinearly gives the three terms.
\end{proof}

This is the floor against which every mask statement in this paper is
judged. It replaces the count-based interval, which is wrong here for
the reason given in Proposition~\ref{prop:coh}.

\subsection{The floor's uncertainty does not fall like a count}

\begin{proposition}[coherent sums]\label{prop:coh}
The error bar of Lemma~\ref{lem:cellmom} does \emph{not} decay like
$n_c^{-1/2}$. For fixed $v$, about $n_c/\log N$ of the terms of
$u_c(v)$ are nonzero, each of size $\log N/\sqrt V$, so
$u_c(v)\approx n_c/\sqrt V$ and
\[
  \frac{Q_{cc}}{n_c^2}
  \;\approx\; \frac{\sum_v\mu^2(v)}{V}
  \;\approx\; \frac{(6/\pi^2)N}{\AAA(N)\,N\log N}
  \;\asymp\; \frac{1}{\log N},
\]
so the standard error falls like $(\log N)^{-1/2}$.
\end{proposition}

Measured over $10^5 \le N \le 1.4\cdot10^7$, fitting
$\mathrm{se}\propto N^{-b}$ gives $b = 0.1105, 0.0465, 0.0385,
0.0379, 0.0378, 0.0379$ at depths $5,4,3,2,1,0$. The law
$(\log N)^{-1/2}$ predicts an \emph{apparent} power exponent of
$1/(2\langle\log N\rangle) = 0.036$ over that range, against the
measured $0.0379$ at the three shallowest depths: agreement to five
percent.

\begin{remark}[a count is not an error bar]
The consequence is not confined to this paper. Over the factor $140$ in
$N$ measured here, $n_c^{-1/2}$ says the error bar shrinks by $11.8$;
it shrinks by $1.21$. \textbf{An interval built from a count is about
ten times too narrow at the top of this range, and the shortfall grows
like $N^{0.46}$.} The cause is that $u_c(v)$ is a coherent sum, not a
self-averaging one, and the same is true of any cell mean of a field
whose summands share a common arithmetic input.
\end{remark}

\subsection{Shape, cause, and size}

\paragraph{Shape.} The cell floor's profile across cells is the
singular series of the shift: correlation between the predicted and
observed profiles is $0.9997$--$1.0000$ in every band measured.

\paragraph{Cause.} The floor collapses to the independent-sign value
$(k-1)/n$ under the permutation of Lemma~\ref{lem:placebo}. So the
floor is carried by the correspondence between cells and divisibility,
not by the cell sizes.

\paragraph{Size.} The floor's magnitude is predicted by
\[
  \delta_c \;\approx\; a\,\bigl(E_{\mathrm{same},c}[\SS_2]
  - E_{\mathrm{all}}[\SS_2]\bigr),
\]
the excess of the singular series of the shift over same-cell pairs,
with $a$ the coefficient in $\rho(h) = a\,\SS_2(h)+b$. Weighting the
cells by $n_c/n$, as the closed form of Lemma~\ref{lem:cellmom}
requires, this accounts for the measured floor to within about six
percent at the top of the range. (A cell-size cap in the sampling ---
weighting cell $c$ by $\min(4n_c, 40000)$ rather than by $n_c$ ---
under-weights the large shallow cells, whose same-cell excess is
smallest, and inflates the prediction by $1.19$ rising to $1.69$ across
the bands. That is the whole of the discrepancy this prediction
originally showed.)

\paragraph{Rarity.} The mask is \emph{rare}, not small. Its per-cell
share reaches $0.94$ in the deepest cell at $N\approx1.4\cdot10^7$,
against a pooled share of $0.018$. Pooling across cells is what makes
it look negligible.

\subsection{The decay, and what does not explain it}

Fitting $|\delta_c| \propto N^{-a_d}$ per depth, with the exact errors
of Lemma~\ref{lem:cellmom} as weights and each exponent's own standard
error from the fit's covariance:

\begin{center}
\begin{tabular}{cccc}
\hline
depth $d$ & $a_d$ & s.e. & $a_d/\mathrm{s.e.}$\\
\hline
5 & $0.1434$ & $0.0155$ & $9.2$\\
4 & $0.2152$ & $0.0065$ & $33.2$\\
3 & $0.2713$ & $0.0040$ & $67.2$\\
2 & $0.3686$ & $0.0052$ & $70.6$\\
1 & $0.0437$ & $0.0556$ & $0.8$\\
0 & $0.6289$ & $0.0121$ & $52.0$\\
\hline
\end{tabular}
\end{center}

A common exponent is rejected at $\chi^2/\mathrm{dof} = 251$ on
$5$ degrees of freedom. Excluding depth $1$, whose exponent is not
measurable, the exponent rises \emph{monotonically as the cell gets
shallower}, with the steps at $5$ to $30$ standard errors: the mask
decays faster where fewer small primes divide $N$.

Two limitations are stated rather than hidden. First, over a factor
$160$ in $N$ the data do not separate $N^{-a}$ from $(\log N)^{-b}$, so
these are exponents \emph{of an assumed form}. Second, the mechanism
that explains the floor's size does not explain its decay:

\begin{proposition}[the size mechanism is scale-invariant]
\label{prop:scaleinv}
$D_c := E_{\mathrm{same},c}[\SS_2] - E_{\mathrm{all}}[\SS_2]$ depends
only on the distribution of the shift $h$ in the residue classes mod
$3,5,7,11,13$ that the cell fixes. That distribution is the same at
every scale, so $D_c$ is scale-invariant and predicts an exponent of
zero at every depth.
\end{proposition}

Computed exactly --- by an FFT autocorrelation of the cell indicator
against $\SS_2$, with no sampling --- the predicted exponents are
$-0.0052, -0.0070, -0.0017, -0.0005, -0.0002, -0.0003$ with standard
errors $0.0001$--$0.0007$, i.e.\ zero at every depth, against measured
exponents of $0.14$ to $0.63$; the residual is $99\%$ of the spread.
\textbf{The mask's decay is therefore not carried by the arithmetic
excess that carries its size.} What produces different exponents at
different depths is not identified.

None of this threatens $C(N)=o(N)$: the mask is lower order under every
parameterisation considered.

\section{The negative map}\label{sec:closures}

Every entry below was pre-registered: the decision rule was fixed in
writing before the computation ran. Route adjudications were done in
fresh context against the source papers' verbatim lemma hypotheses.

\subsection{Adjudication of existing machinery (5)}

\begin{longtable}{p{0.24\textwidth}p{0.10\textwidth}p{0.56\textwidth}}
\hline
Route & Verdict & Blocking coordinate\\
\hline
\endhead
MRT / Lichtman 2020, shift $\to$ dilate & Blocked &
 the orthogonality factorization needs a \emph{linear} pair
 constraint ($h = m-n$); the dilate constraint
 $m'u - mu' = N(m'-m)$ is bilinear. The $h$-average is a translation;
 the $k$-average is a dilation, with no diagonalizing character
 family.\\
Tao 2016 entropy decrement, $k$-averaged & Blocked $\times 3$ &
 (i) no $k$-analog of approximate affine invariance; (ii) the sampling
 prime enters the phase multiplicatively, so the sparsification fails
 and the surviving input is the target itself; (iii) the saving is
 triple-log, below spec.\\
Lichtman 2020 rerun in dilate coordinates & Blocked &
 the congruence coupling is a genuine isomorphism but powers only the
 typical-set restriction; the decoupling blocks at the same bilinear
 coordinate, and the uniformity slot would need an $EH_\mu$-grade
 input (circular).\\
Dirichlet-polynomial fourth moment / Perron & Blocked &
 every log-saving mean-value pillar assumes coefficient
 multiplicativity in its own variable; $\mu(N-u)$ has none. Unfolding
 by Perron costs $T \gtrsim N^{1-o(1)}$; opening the fourth moment
 reproduces the binary correlation.\\
Partial slices & Partial &
 the type-I slice is already consumed; the \emph{$N$-averaged}
 theorem is provable but lands in exceptional-set territory, which
 Huang--Li cannot consume; no nontrivial fixed-$N$ slice exists.\\
\hline
\end{longtable}

\noindent
\textbf{The common obstruction.} $\mu(m)\mu(N-mk)$ couples its
variables simultaneously through the product $mk$ and the difference
$N-mk$; the pair constraint is bilinear and is diagonalized by no
single character family, additive or multiplicative, while each
route's decisive lemma consumes exactly that diagonalization as a
hypothesis.

\subsection{Technique designs, kill-tested (9)}

\begin{longtable}{p{0.05\textwidth}p{0.30\textwidth}p{0.55\textwidth}}
\hline
\# & Design & Result (pre-registered rule applied)\\
\hline
\endhead
K1 & multiplicative Fej\'er kernel on the exact dilation ladder &
 \textbf{dead}: full 63-divisor orbit $R^2 = 0.466$. Half the field's
 energy is invisible to its entire multiplicative orbit.\\
K2 & manufactured pair congruence (determinant/Kloosterman) &
 \textbf{dead}: 0/10 flags over $h = 1\ldots210$; congruent pairs are
 statistically $h$-blind. Its detection floor is
 $4/\sqrt n \approx 0.163$ standard deviations of its field, and the
 design would need $14{,}000$ times that floor at $A=1$ and
 $4.1\cdot10^6$ times at $A=2$, so the verdict is quantitative and not
 merely an absence of flags.\\
K3 & Wishart / operator moment method &
 \textbf{dead}: $\mathrm{tr}(M^2),\mathrm{tr}(M^3),\mathrm{tr}(M^4)$
 dead-center on the Wishart null ($z = -0.48/+0.42/-0.03$). No
 sub-Wishart surplus to fund the exchange.\\
K4 & $N$-average descent &
 \textbf{dead}: 0/8 flags; the dual field is $\delta$-blind and
 $m$-uncorrelated.\\
R1 & zero-spectrum visibility (explicit formula) &
 \textbf{dead}: $R^2_{\text{zeros}} = 0.2152$ vs random-frequency
 $0.2196\pm0.0055$. In standard errors of its own null the measurement
 sits $-0.80$ below, and the data exclude a $7.5\%$ enhancement at
 three standard errors.\\
R2 & determinant / Kloosterman phase &
 \textbf{dead}: regression $R^2 = -0.0001/+0.0004$ against controls
 $\pm0.0002$ over 48{,}000 coprime pairs. The verdict rests on the
 measurement ($-0.38$ s.e.), not on the criterion.\\
R3 & character transform of the $k$-average &
 \textbf{dead by analysis}: for $K \le N^{1/3}$ the transform deposits
 every component into thin progressions (moduli $\ge N^{2/3}$);
 Parseval forbids a statistical gain.\\
R4 & divisor switch &
 \textbf{dead}: see \S\ref{sec:R4}.\\
R5 & the circle method applied directly to $C(N)$ &
 \textbf{dead, zero margin}: both bills lie at or above the trivial
 bound, and the cap on the pointwise route is Parseval
 ($\sup\ge\|\cdot\|_2$), which no technique can move
 (Proposition~\ref{prop:E}).\\
\hline
\end{longtable}

\begin{remark}[what a null verdict costs]
A kill-test that does not fire is evidence of absence only against a
stated floor. Where a design's threshold was chosen as an effect size
rather than in standard errors of its own null, we restate it: R4 is a
five-sigma test at block size $B=8$ and carries no information at
$B=512$, its precision degrading as $\sqrt{2/n_b}$; and the C4
threshold of $0.5\times$ sits only $1.29$ standard errors below its
null, so noise satisfies it $8.8\%$ of the time. In both cases the
DEAD direction stands on the measurement --- never systematically below
the null --- and not on the count of flagged levels.
\end{remark}

\subsection{Representation classes (3, plus one open)}

\begin{longtable}{p{0.10\textwidth}p{0.25\textwidth}p{0.55\textwidth}}
\hline
Class & Test & Result\\
\hline
\endhead
C-I abelian & rational-peak energy vs mask-null &
 \textbf{closed}, 0/6: the abelian spectrum is mask-exact.\\
C-II inverse domain & special-frequency excess in the
 modular-inverse re-indexing &
 \textbf{closed on verification}: fired at 4/300 with $z = +10.97$
 under 8-draw nulls, collapsed under 64-draw nulls
 ($+4.14/+2.41/+1.19/+0.30$), permutation null concurring, second-$N$
 replication 0/4.\\
C-IV manufactured modularity & approximate Fricke law for
 $\Phi(z) = \sum t_m e(mz)$ &
 \textbf{closed}, 0/6 levels. A true modular form drives the defect to
 $\approx 0$, so absence is meaningful.\\
C-III Motohashi type & (no finite test) &
 \textbf{open}, needing exactly the three items of \S\ref{sec:c3}.\\
\hline
\end{longtable}

\subsection{R4 in detail: the divisor switch does not
localize}\label{sec:R4}

Applied to the dilate field over its \emph{full} ranges the switch
gives an exact identity
\[
  \sum_{k\ge1}\ \sum_{m:\,mk\le N-1}\mu(m)\mu(N-mk)
  = \sum_{u<N}\mu(N-u)\sum_{m\mid u}\mu(m)
  = \mu(N-1),
\]
verified by brute force at $N = 5000$ and $N = 20000$. This is perfect
cancellation: $O(1)$ for a double sum of $\sim N\log N$ terms, far
beyond square root, and it is the strongest cancellation found
anywhere in this work.

E1 imposes two restrictions that make the inner divisor sum
incomplete: the type-II cut $m > \sqrt N$ and the dyadic band
$k \sim K$. The kill-test asked whether any of the cancellation
survives, by measuring block sums $S_B(j) = \sum_{k\in\text{block}}D(k)$
against the $B$-independence that Conjecture~\ref{conj:L} predicts.

\emph{Result: none of it survives.} On $8000$ values of $k$ --- the
entire band on which the type-II field is non-empty, since
$m > \sqrt N$ forces $k < \sqrt N$ --- the block ratios are flat
($B=8$: $0.958$ and $1.023$ at two $N$, against $B=1$ baselines
$0.980$ and $0.979$), and the sharper diagnostic, the lag-1
autocorrelation of $D(k)/\sqrt{\mathrm{supp}(k)}$, reads
$+0.0104$ and $+0.0127$ against a standard error of $0.0112$: dead
zero, and if anything mildly positive rather than the negative
coherence a surviving residue would require.

\emph{The mirror.} Switching the banded $L^1$ sum gives
$\sum_{k\sim K}D_{\text{full}}(k)
 = \sum_{u<N}\mu(N-u)\sum_{m\mid u,\ u/2K<m\le u/K}\mu(m)$, where
$m \asymp N/K \ge N^{2/3}$: the surviving M\"obius sits on the
\emph{long} variable, the exact opposite of the assignment that makes
Theorem~\ref{thm:A} work. The switch is not a technique that happens
to fail on the supply side; its single requirement --- M\"obius on the
short variable --- is precisely what the type-II cut forbids.

\subsection{C-III: what it needs}\label{sec:c3}

C-III is the one representation class with no finite test, and it is
open. Three requirements are identified; two are settled here against
the natural construction, and the third is external.

\paragraph{(1) A legitimate transform --- closed for changes of
variable.} Put $A = N-a$, $B = N-b$, so $A = mk$ and $B = m'k$; then
$m'A - mB = 0$ exactly. In centered coordinates the dilate family is
therefore the \emph{pencil of lines through the origin}, while the
shift family $B = A - h$ is a family of \emph{parallel} lines. A
pencil is characterised by its vertex, here a finite point; a parallel
family is a pencil whose vertex is at infinity. Affine maps send
finite points to finite points, and since $\mu$ lives on $\mathbb Z$
the available changes of variable are exactly the integral affine
ones. So no change of variable carries one family to the other. This
proves what a direct probe had only measured: the obvious re-indexing
is circular, the off-diagonal computed directly and via that
re-indexing agreeing exactly. What remains is the summation-formula
class, blocked in general by the roughness of the outer weight $\mu$,
with every named candidate in it already closed.

\paragraph{(2) A classification covering the type-II region ---
settled against the natural bookkeeping.} That bookkeeping needs the
M\"obius-side $a$ to satisfy $a \le y^{O(1)} = M^{o(1)}$. Measuring the
absolute Heath--Brown weight
$W(a) = \sum_j \binom{J}{j} A_j(a) D_j(M/a)$ across $a$-sizes, the
fraction with $a > M^{0.05}$ is $0.939, 0.949, 0.960, 0.947$ at
$J = 3, 4, 6, 8$, rising to $0.961$ and $0.969$ at $M = 10^6$. So
about $95\%$ of the identity's weight lies outside the region the
classification covers, flat in $J$ across the admissible range and
increasing with $M$. Those pieces carry a rough coefficient with no
divisor structure and hence no Voronoi entry.

The obstruction is parameter-independent: the identity with cut $z$
needs $z^J \ge x$, and its $j$-th term has $a \le z^j$, which at
$j = J$ is $x$ for every admissible $(z,J)$ --- while the weight
concentrates in exactly those high-$j$ terms ($0.824$ at $j=3$ when
$J=3$; $0.847$ at $j \in \{6,7,8\}$ when $J=8$). A repair would have
to bound $\sum_a\sum_b \alpha(a)\beta(b)\mu(N-abk)$ with $\alpha$
rough and $a$ long, where every individual piece is trivial and all
the content is cancellation across $a$ --- a type-II estimate for
$\mu(N-\cdot)$, i.e.\ the wall. Every identity decomposing $\mu$
produces such a term; one whose every piece had either a long free
variable or a divisor-structured rough coefficient would dispose of
the parity obstruction. \textbf{Completing the construction and
breaking the wall are the same task.}

\paragraph{(3) Quantitative averaged Chowla.} Shift-averaged binary
$\mu\mu$ correlations at a \emph{fixed} log-power saving, against a
best known of $(\log)^{1-c}$. This is a named external open problem
and is where C-III stands.

That third item sets the character of the obstruction. The
adjudication's central finding was that the \textbf{dilate} average
admits no diagonalizing character family; but the \textbf{shift}
average is precisely the home ground of the
Matom\"aki--Radziwi\l\l--Tao machinery. If a legitimate transform
converting one into the other exists, what remains is quantitative
strength inside an active area rather than the absence of any coupling
surface.

\section{What the closures constrain}

Any transform or invariance $T$ that linearizes the pair constraint at
a cost of at most a log power must satisfy all five of the following.

\begin{enumerate}
\item \textbf{(K1)} $T$ cannot act through divisibility sub-sums: the
  multiplicative orbit spans only half the field's energy, so $T$ must
  couple to the $p$-coprime core directly.
\item \textbf{(K2)} $T$ cannot manufacture its congruence externally:
  collapse structure pays only when it arises inside an intrinsic
  average that is already present.
\item \textbf{(K3)} $T$ cannot bootstrap from the field's own moments:
  every moment consumes higher $\mu$-correlations and the field
  carries no sub-Wishart surplus.
\item \textbf{(K4)} $T$ cannot descend from the $N$-average: its
  linearization is load-bearing and the $k$-average holds no shadow
  of it.
\item \textbf{(R4)} $T$ cannot be the divisor switch or a relative:
  the switch's cancellation is a property of \emph{complete} divisor
  sums, and the type-II cut makes every divisor sum incomplete.
\end{enumerate}

\noindent
Together with the round-2 findings --- zeros invisible through the
pairing, characters merely relocating the difficulty into thin
progressions, determinant phases blind --- the field offers no
coupling surface in any direction that has an existing mathematical
name. In the automorphic world, shifted convolutions
$\sum a(n)a(n+h)$ are controlled because the coefficients $a(n)$
\emph{come with} a spectral representation; $\mu$ has none, and every
route above failed exactly where it tried to borrow one. The technique
to be created is a spectral (or equivalent structural) representation
of $\mu$-pairs themselves, not a projection onto existing spectra ---
and its construction is a purely theoretical act, beyond what
kill-testing can reach.

\section{The margin, and where the difficulty is}

The requirement constrains every $N$, so the figure to quote is the one
at the extreme, not at a typical $N$. With $\max|G|\approx a_n$ the
Gumbel location,
\[
  \max_{N\le X}|C(N)| \;\approx\; a_n\sqrt{\AAA(N)\,N\log N},
  \qquad
  \frac{N}{\max|C|}\;\approx\;
  \frac{\sqrt N}{a_n\sqrt{\AAA\log N}},
\]
which is $10^{4.4}$ at $N=10^{12}$ and $10^{22.8}$ at $N=10^{50}$. The
requirement is not remotely tight; measured, $\max|C|/N$ falls from
$0.056$ to $0.0082$ over the range computed, and the margin at
$N=10^8$ is a factor $N^{0.454}$.

Placing the two standard estimates in the same units makes the shape of
the difficulty explicit. The target is $N(\log N)^{-A}$. The trivial
bound sits a factor $(\log N)^{A}$ above it; Cauchy--Schwarz gives
$N\sqrt{6(\log N-1)/\pi^2}$, a factor $(\log N)^{A+1/2}$ above it; and
the truth sits a \emph{power of $N$ below} it. \textbf{The entire
difficulty is a power of $\log N$}, and it is not the size but the
proof that is missing --- by Theorems~\ref{thm:D} and~\ref{thm:Dprime}
the one classical route to proving it is closed over its whole design
space.

\begin{remark}
Everything measured here is at $N\le1.6\cdot10^7$, with two arms at
$N\approx10^8$. The no-go theorems of Section~\ref{sec:demand} are
asymptotic and acquire content only near $N\approx10^{480}$; nothing
measured here constrains that range, and nothing there is contradicted
by these measurements.
\end{remark}

\section{Methodology}

Four rules were followed throughout, and they are the reason we believe
the negative results.

\begin{enumerate}
\item \textbf{Pre-registration.} Every design's decision rule was
  written before its computation ran, including the rule that would
  refute the hypothesis under test. Two designs fired and were then
  retracted on verification (C-II at $z = +10.97$; a $z = 9$ spectral
  excess that was a null-design category error --- an i.i.d.-entry
  Wigner null applied to a Gram matrix).
\item \textbf{Adversarial review in fresh context}, against the source
  papers' verbatim lemma hypotheses rather than against summaries of
  them.
\item \textbf{Power before belief.} A $+3.7\sigma$ band elevation
  regressed to $1.078\pm0.072$ under bootstrap; a $-2.4\sigma$ dip
  resolved at four times the sample; the R4 deficit dissolved at four
  times the $k$-range.
\item \textbf{Nulls before thresholds, and weights before
  comparisons.} A threshold means nothing until the spread of the
  quantity it judges has been measured (Remark~\ref{rem:scale} and the
  restatements in \S\ref{sec:closures}); and two summaries of one
  object are not comparable until each one's weight is stated
  (Remark~\ref{rem:rho}). Both rules were adopted after the
  corresponding mistakes were made and caught.
\end{enumerate}

\noindent
Two traps are worth recording because they are traps and not slips.

\begin{itemize}
\item \emph{The $(q,N)>1$ main-term trap.} In the proof of
  Theorem~\ref{thm:A}, assigning a main term $T_m/\varphi(q)$ to a
  residue class with $(q,N)>1$ --- which is degenerate, its true
  contribution being $O(\log^2 N)$ --- shifts the density by exactly
  $N/\varphi(N)$. Carried into the $\log k$ branch, that error
  produces an apparent \emph{refutation of $EH_\mu$}. This is the most
  plausible false positive we met.
\item \emph{Eight-draw-null inflation.} Nulls estimated from too few
  draws inflate maxima across a family of 300 tests; the C-II ``hit''
  survived 8-draw nulls and died under 64-draw nulls plus replication.
\end{itemize}

\section{Reproducibility}

All measurements run on a laptop with Python and numpy. The
repository's \texttt{PROVENANCE.md} maps every numbered statement in
this paper to the code that verifies it and the result file it was
read from. One-shot verification of the core corpus is
\texttt{python v1/code/verify/verify\_all.py}, and the deep-$N$ arm is
\texttt{python v1/code/verify/verify\_deep.py}. The consistency gate, which runs
every consistency checker and exits nonzero if any fails, is
\texttt{python v1\_log/code/gate.py}, and belongs to the program's
own record rather than to this paper.

\section{Summary}

\begin{itemize}
\item The demand side of the Huang--Li reduction is closed: one half
  is now an unconditional theorem, the other half is equivalent to the
  conclusion.
\item The supply side is closed to measurement: the object is
  featureless in every direction with a name, and the eighteen
  pre-registered closures say precisely which directions were checked
  and why each fails.
\item The wall has an exact second moment, an exact aggregate identity,
  a Gaussian bulk and tail under that scale and no other, and a
  deterministic mask whose shape, cause and size are accounted for and
  whose decay law is not.
\item Net progress toward Goldbach is zero. The contribution is the
  map, the law, one unconditional theorem that removes a
  Goldbach-neutral demand, one conjecture that any future construction
  must reproduce, and one defect report for the source paper.
\end{itemize}

\begin{thebibliography}{9}
\bibitem{HL} Jing-Jing Huang and Huixi Li, \emph{On the connection
between the Goldbach conjecture and the Elliott--Halberstam
conjecture}, arXiv:2005.03811v2 [math.NT], 2022.
\bibitem{ThmA} \emph{An unconditional bound for a M\"obius-weighted
correlation sum in fixed residue classes}, companion note,
\texttt{paper/v1/theorem\_A.tex} in this repository.
\bibitem{Pan} Cheng-Dong Pan, \emph{A new attempt on Goldbach
conjecture}, Chinese Ann. Math. \textbf{3} (1982), 555--560.
\bibitem{Tao} Terence Tao, \emph{The logarithmically averaged Chowla
and Elliott conjectures for two-point correlations}, Forum Math. Pi
\textbf{4} (2016).
\bibitem{Li20} Jared Duker Lichtman, \emph{Averages of the M\"obius
function on shifted primes}, arXiv:2009.08969.
\bibitem{Mir49} L. Mirsky, \emph{The number of representations
of an integer as the sum of a prime and a $k$-th power free
integer}, Amer. Math. Monthly \textbf{56} (1949), 17--19.
\bibitem{MRT15} Kaisa Matom\"aki, Maksym Radziwi\l{}\l{}
and Terence Tao, \emph{An averaged form of Chowla's conjecture},
Algebra Number Theory \textbf{9} (2015), 2167--2196.
\bibitem{Li23} Jared Duker Lichtman, \emph{Primes in arithmetic
progressions to large moduli, and shifted primes without large prime
factors}, arXiv:2309.08522.
\bibitem{MV17} M. Ram Murty and Akshaa Vatwani, \emph{Twin primes and
the parity problem}, J. Number Theory \textbf{180} (2017), 643--659.
\end{thebibliography}

\end{document}
